% 第22段階の追加記録. 第21段階の本文に続けて追加する.
% 既存の \bm, \bbo, booktabs, hyperrefを使用する.
% 新しい引用キー, 文献環境, 文書終了命令は追加しない.
% 根拠資料: bivariate_survival_reconstruction_stage22_resumed.zip.
% 以下の数値は第22段階とその再開時の保存結果である.

\section{第22段階の目的と連続母標本による統合検証}
\label{sec:stage22_scope}

本段階では, 連続母分布から個々の潜在天体を生成し,
非入れ子の検出限界, A--D選択, 区間記録,
学習予測, 親質量の点推定と親信頼集合への真値所属を接続する.
全120標本で点推定と真値所属を検査し,
事前指定した3標本では保証付き生存射影まで計算した.
観測法則への適合と親分布全体の回復を別々に評価する.

検出条件は$Z_b\le C_b$とし, 等号を検出に含める.
A, B, Cだけを収録し, Dの個数と潜在値は推定器へ渡さない.
親変数と限界の独立性を維持し,
検出値は事前指定された区間番号として記録する.
未群化測定値の密度尤度や測光誤差を扱う実験ではない.

\section{観測設計と推定モデル}
\label{sec:stage22_design}

両座標の記録区間を$(-\infty,1],(1,\infty)$とし,
限界と対象点を
\begin{align}
\bm c_1&=(2,1),\qquad \bm c_2=(1,2),\notag\\
\bm c_3&=(2,2),
\label{eq:stage22_limits}\\
\bm t_1&=(1,1),\qquad \bm t_2=(3/2,1),\notag\\
\bm t_3&=(2,2)
\label{eq:stage22_targets}
\end{align}
とする. 最初の二つの収録集合は互いに包含されない.
全可能非D記録は限界別に5,5,8種類である.
各記録の整合集合$E_{jr}$は$D_j^c$を分割し,
\begin{align}
\pi_P(r\mid\bm c_j)&=\frac{P(E_{jr})}{P(D_j^c)},\qquad
\sum_r\pi_P(r\mid\bm c_j)=1
\label{eq:stage22_record_law}
\end{align}
を満たす. 未出現記録も正規化と候補モデルに保持する.

限界, 記録境界と対象点から全実数空間を35セルへ分割する.
全観測列を保存する同値類は9個,
三対象の生存指示も保存する同値類は11個である.
負座標, 境界単点, 非有界領域と完全不可視領域も残す.
真の支持に合わせた切り詰めや, 優越性による射影用支持の削除は行わない.
親信頼クラスの外部条件と有意水準は
\begin{align}
P(D_j)&\le\frac12\quad(j=1,2,3),\qquad
\alpha_{\mathrm{sig}}=\frac1{20}
\label{eq:stage22_tail_class}
\end{align}
とする. 尾部上界はデータから得た収録率ではない.

\section{評価真値と選択前の標本生成}
\label{sec:stage22_generation}

評価専用の連続密度を
\begin{align}
f_\lambda(z_1,z_2)
&=\frac1{16}\left[1+\lambda
\left(1-\frac{z_1}{2}\right)\left(1-\frac{z_2}{2}\right)\right]
\bbo\{(z_1,z_2)\in(0,4)^2\},
\label{eq:stage22_true_density}\\
\lambda&\in\{-3/4,0,3/4\}
\label{eq:stage22_lambda}
\end{align}
とする. 密度族と支持$(0,4)^2$は真値評価にだけ用いる.
真のセル質量は矩形積分で有理数として求める.
真の不可視確率は
\begin{align}
P_0(D_1)=P_0(D_2)&=\frac38+\frac{3\lambda}{64},\qquad
P_0(D_3)=\frac14+\frac{\lambda}{16}
\label{eq:stage22_true_missing}
\end{align}
であり, 三真値とも宣言した親クラスに属する.

独立一様乱数$U,W$を使い,
\begin{align}
a&=\lambda(1-2U),\qquad
V=\frac{2W}{1+a+\sqrt{(1+a)^2-4aW}},
\label{eq:stage22_sampler}\\
\bm Z&=(4U,4V),\qquad V+aV(1-V)=W
\label{eq:stage22_conditional_cdf}
\end{align}
によって生成する. 生成は倍精度であり,
保存値の条件付きCDF恒等式の最大残差は約$3.33\times10^{-16}$だった.

親標本数を400,1600とし, 各$\lambda$・各標本数で20反復を行う.
合計120標本, 120000親天体である.
次の種から三つの乱数生成器を独立に分岐する.
\begin{verbatim}
SeedSequence([20260913, lambda_index, N_parent, rep])
\end{verbatim}
限界を三種類から等確率で割り当て,
学習所属を確率$1/2$で割り当てる. 両者とも潜在値と独立である.
学習・検証の分割は選択前に行い, その後Dを除外する.
従って収録数, 限界別件数と学習・検証件数はランダムであり,
収録数固定の多項個数実験とは異なる.

\section{学習予測と点推定の役割分離}
\label{sec:stage22_estimation}

学習予測は観測9クラスの一様初期質量から,
学習記録だけに従来のD補完EMを有理数で2回適用する.
全データ点推定には学習・検証個数の和を使い,
同じ原EMを平均正スコア$10^{-8}$, 最大4000更新で実行する.
全データ適合を検証尤度比の予測側へ戻さない.

原EMは尾部制約付き最適化ではない.
最終候補の尾部所属を別途点検し, 今回は120候補全てが条件を満たした.
全標本で25--53更新の間に数値停止したが,
スコア停止を大域MLEの証拠へ読み替えない.

参考生存曲面には, 観測クラスの総質量を保存して
対象保存クラスへ等配分する事前規則を用いる.
対象$(3/2,1)$の内部配分は観測から推定されたものではない.
真の記録法則への層件数重み付きTVと,
三固定点$\mathcal T=\{\bm t_1,\bm t_2,\bm t_3\}$での参考生存RMSEを別に記録する.
選択後の隠れた潜在値による経験分布と,
Dを含む全親標本の経験分布も評価専用に区別する.

\section{親信頼集合の真値所属と標本結果}
\label{sec:stage22_coverage}

学習予測を$\widetilde P$として, 従来と同じ
\begin{align}
\mathcal C
&=\left\{P:P(D_j)\le\frac12\ \text{for all }j,
\ L_{\mathrm{val}}(P)\ge
\frac1{20}L_{\mathrm{val}}(\widetilde P)\right\}
\label{eq:stage22_confidence_set}
\end{align}
を用いる. 全記録法則の正規化と独立な学習・検証分割が校正の前提である.
真の$P_0$の所属は
\begin{align}
\log E_0
&=\sum_{j,r:N_{jr}>0}N_{jr}
\log\frac{\pi_{\widetilde P}(r\mid\bm c_j)}
{\pi_{P_0}(r\mid\bm c_j)}
\label{eq:stage22_true_membership}
\end{align}
を有理数区間で囲み, $\log20$と比較して判定する.
予測尤度がゼロなら棄却しないという従来の規約を保持する.

\begin{table}[htbp]
\centering
\caption{第22段階の20反復平均と親信頼集合への真値所属.
RMSEは三点の参考曲面についての値である.}
\label{tab:stage22_sampling}
\begin{tabular}{rrrrr}
\toprule
$\lambda$ & 親標本数 & TV誤差 & 参考RMSE & 真値採択 \\
\midrule
$-3/4$ & 400  & 0.059385 & 0.076564 & 20/20 \\
$-3/4$ & 1600 & 0.029562 & 0.073520 & 20/20 \\
$0$    & 400  & 0.063484 & 0.111616 & 20/20 \\
$0$    & 1600 & 0.033199 & 0.108534 & 20/20 \\
$3/4$  & 400  & 0.064485 & 0.148211 & 20/20 \\
$3/4$  & 1600 & 0.027549 & 0.146064 & 20/20 \\
\bottomrule
\end{tabular}
\end{table}
全120標本で真値を採択し, 算術的な所属判定未確定はなかった.
これは各設定20反復の実現値であり, 被覆確率1という結論ではない.
全個数表による正確な被覆列挙でもない.

\section{不可視質量による参考推定誤差の下限}
\label{sec:stage22_nonidentification}

全限界で完全不可視な領域$I=(2,\infty)^2$について,
原EMは初期質量を保存するため,
\begin{align}
\widehat P(I)&=\frac19,\qquad \widehat S(2,2)=\frac19
\label{eq:stage22_invisible_initialization}
\end{align}
となる. 真の値は$\lambda=-3/4,0,3/4$の順に
$13/64,1/4,19/64$である.
この相違は標本数や反復数を増やすだけでは消えない.
三点の参考RMSEには
\begin{align}
\operatorname{RMSE}_{\mathcal T}(\widehat S,S_0)
&\ge\frac{|S_0(2,2)-1/9|}{\sqrt3}
\label{eq:stage22_rmse_floor}
\end{align}
という下限があり, 順に約0.053124,0.080188,0.107251となる.
これは今回の初期分布と三点評価に対する代数的な下限である.

観測法則へのTV誤差が減っても, 親分布の参考RMSEには
不可視質量と記録内部配分の自由度が残り得る.
一方, 同じ条件付き観測法則を許す親信頼集合は真値を含み得る.
一つの代表分布と, 許される分布集合を区別する必要がある.

\section{事前指定した三標本の保証付き生存射影}
\label{sec:stage22_projection}

各$\lambda$の親400天体・反復番号0を事前指定し,
対象$S_P(1,1)$を射影する.
学習予測, 尾部条件と閾値は各標本のものを保持する.
第20段階の無改変エンジンは15ノードで三例とも精度未達成だった.
その証拠を祖先として, 第21段階の無改変エンジンを相対尺度で用い,
各端点幅$10^{-3}$, 最大65ノード, 提案反復上限150で実行した.

\begin{table}[htbp]
\centering
\caption{親生存信頼射影の端点位置の囲い.
小数は外向き8桁表示であり, 区間全体の幅と端点計算誤差を区別する.}
\label{tab:stage22_projection}
\small
\begin{tabular}{rllr}
\toprule
$\lambda$ & 下端の囲い & 上端の囲い & ノード \\
\midrule
$-3/4$ & $[0.17043061,0.17100683]$ & $[0.70439875,0.70493109]$ & 23 \\
$0$ & $[0.21144468,0.21223209]$ & $[0.75014270,0.75093018]$ & 13 \\
$3/4$ & $[0.20586024,0.20658897]$ & $[0.69722377,0.69792403]$ & 15 \\
\bottomrule
\end{tabular}
\end{table}
三例とも指定精度に達し, 外側範囲は真値を含んだ.
真の$S_0(1,1)$は順に$549/1024,9/16,603/1024$である.
他の117標本の鋭い射影は計算しておらず,
三例の結果を120例の射影被覆率へ拡張しない.
生存範囲自体に残る広さを, 端点計算の未収束とは解釈しない.

\section{独立検証, 再開時の再実行と保存範囲}
\label{sec:stage22_record}

保存した潜在配列から120公開入力を再構成した.
非検出成分とDの潜在値を増加させても公開入力が変わらないことを120件確認し,
2880点の記録写像と54記録確率の直接積分を照合した.
不可視初期質量の最大変動は約$6.245\times10^{-16}$だった.
生の最小尤度増分$-4.5475\times10^{-13}$は丸め診断として保持し,
履歴を単調化していない.

三射影の独立検証は, 祖先証拠を含めて
51ノード,54双対証拠,594列不等式,82接線と6内側親分布を確認した.
検証中はLP, 非線型最適化と頂点列挙を禁止する.
学習2EMの定義の再計算と有理数演算は行う.

当初の記録では標本計算を再実行し, 射影は証拠の再検査に留めた.
応答再開時には追加で, 全120標本の生成・予測・点推定・真値所属,
三つの初期射影と三つの改良射影の探索を再実行した.
120入力JSON,120適合JSONと2CSVは全バイトが一致し,
600潜在配列は内容が一致した.
六つの射影証拠と集約JSONは実行時間以外が一致した.
NPZの時刻情報を含む全バイト一致を主張するものではない.

保存パッケージは次のZIPである.
\begin{verbatim}
bivariate_survival_reconstruction_stage22_resumed.zip
\end{verbatim}
\texttt{audit\_stage22.py}に標本実験,
\texttt{project\_stage22.py}と\texttt{refine\_projections\_stage22.py}に射影,
\texttt{verify\_stage22.py}に独立検証を保存した.
元の出力は\texttt{outputs/}, 再開時の出力は
\nolinkurl{resumption_checks/sampling_rerun/}へ分離した.
再実行の一致は
\nolinkurl{resumption_checks/resumption_reproducibility.json}に記録している.
本LaTeX記録では両方の実施範囲を区別して収録し,
その作成だけを新しい標本実験とは数えない.

到達点は, 非入れ子限界と連続母分布について,
選択から点推定・親信頼集合・代表標本の射影までを接続したことである.
原EMの一般的一致性, 多様な母分布での精密な被覆率,
未群化測定値と測光誤差まで検証したものではない.
次段階では同じ潜在標本を用いて観測深度を比較し,
観測情報の増加と, 不可視質量・区間内部の不定性を分けて調べる.
